\documentclass[11pt]{article}
\usepackage[utf8]{inputenc}
\usepackage[T1]{fontenc}
\usepackage{amsmath,amssymb,amsthm}
\usepackage{booktabs}
\usepackage{graphicx}
\usepackage{multirow}
\usepackage{xcolor}
\usepackage[margin=1in]{geometry}
\usepackage{natbib}
\usepackage{hyperref}
\usepackage{tabularx}
\usepackage{array}
\usepackage{adjustbox}

\newtheorem{theorem}{Theorem}
\newtheorem{proposition}{Proposition}
\newtheorem{corollary}{Corollary}
\newtheorem{definition}{Definition}
\newtheorem{conjecture}{Conjecture}

\newcommand{\bC}{\mathbf{C}}
\newcommand{\bI}{\mathbf{I}}
\newcommand{\bB}{\mathbf{B}}

\title{Multi-Class Annotation Errors Are Not Created Equal:\\Bias Taxonomy, Sensitivity Analysis, and Correction Guidance\\for LLM-Annotated Causal Studies}

\author{Anonymous}

\date{}

\begin{document}

\maketitle

%% ============================================================
%% ABSTRACT
%% ============================================================
\begin{abstract}
Large language models (LLMs) are rapidly replacing human annotators in computational social science, yet recent audits show that 31\% of downstream causal conclusions drawn from LLM-annotated data are incorrect.
We show that the classical guarantee---non-differential misclassification biases toward the null---breaks when annotated variables have $K\!\ge\!3$ categories.
The bias behavior is \emph{topology-dependent}: the same confusion matrix that produces innocuous residual bias in a confounding DAG can produce direction-reversing bias in mediation or collider DAGs.
We make three contributions.
\textbf{(1)}~A \emph{multi-class bias taxonomy}: closed-form bias expressions for $K\!=\!3$ and matrix-form expressions with eigendecomposition-based interpretability tools for general $K$, across six canonical DAG topologies.
Four of six topologies (exposure, mediation, collider, front-door) exhibit direction reversal under non-differential $K\!\ge\!3$ error; confounding and instrument topologies preserve monotonicity.
\textbf{(2)}~\emph{Annotation Sensitivity Values} (ASVs), analogous to E-values for unmeasured confounding, that quantify how much annotation error is needed to overturn a conclusion.
We prove monotonicity and sharpness for three priority topologies and provide MC-validated conjectures for the remaining three.
\textbf{(3)}~A \emph{correction guidance framework} mapping \{topology $\times$ error structure $\times$ ASV\} to the appropriate correction method, achieving 95\% recommendation accuracy across 432 simulated scenarios and validated on three CSS tasks annotated by three LLMs (${\sim}270$K annotations).
Binarizing a $K$-class variable to retreat into the binary safe zone is itself hazardous: collapsing $K\!=\!7$ to binary introduces up to $\Delta\!=\!0.19$ additional bias, with 12\% of scenarios exhibiting direction reversal from binarization alone.
We release \texttt{annot-sensitivity}, an open-source Python toolkit implementing the full framework.
\end{abstract}

%% ============================================================
%% 1  INTRODUCTION
%% ============================================================
\section{Introduction}
\label{sec:intro}

LLMs have become the default annotation engine in computational social science (CSS).
Their speed and cost advantages are undeniable, but measurement error in LLM annotations propagates into downstream causal estimates.
\citet{baumann2025} audit this pipeline and find that 31\% of downstream conclusions are already incorrect.
\citet{kim2025} show that cross-model annotation errors are highly correlated---when two LLMs both err, 60\% of the time they err toward the same category---meaning multi-model ensembles cannot eliminate systematic bias.

The classical measurement-error literature provides reassurance for the binary case: non-differential misclassification of a binary exposure biases toward the null \citep{bross1954}.
Practitioners routinely extend this guarantee to multi-class variables by binarizing: collapsing $K$-class annotations into two categories before analysis.
We show that both the original guarantee and the binarization escape route fail for $K\!\ge\!3$.

\paragraph{Core finding.}
Non-differential $K\!\ge\!3$ misclassification breaks the toward-null guarantee in 4 of 6 canonical DAG topologies.
The reversal pattern is strictly topology-dependent: the same confusion matrix that is benign in the confounding topology produces direction-reversing bias in mediation and collider topologies.
In our Monte Carlo experiments, the mediation topology exhibits a 52\% violation rate (47\% sign reversal), the collider topology 41\% (36\% sign reversal), the front-door topology 34\% (28\% sign reversal), and the exposure topology 85\% (83.4\% sign reversal) under 100K random confusion matrices.
Confounding and instrument topologies maintain zero violations.

\paragraph{Why binarization is not a safe retreat.}
Collapsing a $K$-class variable to binary discards inter-category bias structure.
We show that binarization introduces additional bias up to $\Delta\!=\!0.19$ (mediation, $K\!=\!7$), and in 12\% of $K\!=\!5{,}7$ scenarios the binarization itself flips the bias direction---making the collapsed estimate \emph{worse} than the multi-class estimate.

\paragraph{Contributions.}
\textbf{(1)}~A multi-class bias taxonomy with closed-form and matrix-form bias expressions across six DAG topologies (\S\ref{sec:taxonomy}).
\textbf{(2)}~Annotation Sensitivity Values (ASVs) with formal monotonicity and sharpness proofs for three topologies and MC-validated conjectures for three more (\S\ref{sec:asv}).
\textbf{(3)}~A correction guidance framework achieving 95\% recommendation accuracy (\S\ref{sec:correction}).
We validate all results on empirical confusion matrices from three LLMs across three CSS tasks (\S\ref{sec:empirical}) and release the \texttt{annot-sensitivity} Python toolkit.


%% ============================================================
%% 2  RELATED WORK
%% ============================================================
\section{Related Work}
\label{sec:related}

\paragraph{Measurement error in causal inference.}
The epidemiological literature on exposure misclassification dates to \citet{bross1954}, with \citet{brenner1993} first noting that polytomous confounder misclassification can bias in either direction.
\citet{vanderweele2019} provide sensitivity bounds for binary risk ratios.
None of these works provide systematic multi-class bias expressions across multiple DAG topologies.

\paragraph{Annotation error in NLP/CSS.}
\citet{wooddoughty2018} apply DAG reasoning to text classifier errors but produce no formal bias theorems and are limited to binary variables.
\citet{teblunthuis2024} provide the most complete binary treatment with MLA correction but are limited to $K\!=\!2$ and a single classifier.
\citet{lee2024} extend to categorical variables and LLM classifiers but cover only one DAG topology (unobserved confounding) on one dataset, with no systematic taxonomy or sensitivity analysis.
\citet{jerzak2025} benchmark correction methods but treat annotation error as monolithic---if different confusion matrix structures produce qualitatively different bias in different topologies, unconditional benchmarking is itself misleading.

\paragraph{Sensitivity analysis.}
E-values \citep{vanderweele2017} quantify robustness to unmeasured confounding via a single scalar.
Our ASVs serve an analogous role for annotation error but require a reference confusion matrix and amplification direction, reflecting the richer parameterization of $K\!\times\!K$ error structures.


%% ============================================================
%% 3  MULTI-CLASS BIAS TAXONOMY
%% ============================================================
\section{Multi-Class Bias Taxonomy}
\label{sec:taxonomy}

\subsection{Setup and Notation}

Let $\tilde{V}$ denote the observed (LLM-annotated) version of a $K$-class variable $V$.
The misclassification process is characterized by the $K\!\times\!K$ confusion matrix $\bC$, where $C_{jk} = P(\tilde{V}=k \mid V=j)$.
Non-differential misclassification requires $P(\tilde{V} \mid V, \text{other variables}) = P(\tilde{V} \mid V)$, i.e., the error depends only on the true value.
We consider six canonical DAG topologies where $V$ plays distinct structural roles: confounder, treatment (exposure), mediator, collider, instrument, and front-door criterion variable.

For $K\!=\!3$, we derive closed-form bias expressions as functions of $\bC$ and the structural parameters of each DAG.
For general $K$, we provide matrix-form expressions and develop three interpretability tools: (a) eigendecomposition summaries, (b) worst-case scalar bounds, and (c) spectral diagnostic indicators.

\subsection{Closed-Form Bias Expressions (K=3)}
\label{sec:kthree}

For each topology, we derive the bias $B(\bC) = \hat{\theta}(\tilde{V}) - \theta^*(V)$ as a function of $\bC$ and the true structural parameters.
The derivation proceeds by expressing the observed conditional distributions in terms of the true distributions and $\bC$, then computing the resulting distortion in the estimand.

\paragraph{Validation.}
We validate all closed-form expressions via Monte Carlo simulation: 6 confusion matrix structures $\times$ 6 topologies $\times$ 1000 replications at $N\!=\!10{,}000$.
The theory-MC match rate is 97.2\% (35/36 structure$\times$topology combinations fall within 95\% MC confidence intervals).
The single boundary case---front-door with rank-deficient $\bC$---shows a discrepancy of 0.003 above the CI upper bound, which resolves at $N\!=\!50{,}000$.
Table~\ref{tab:theory-mc} reports the full match matrix.

\subsection{Topology-Dependent Reversal}
\label{sec:reversal}

We execute a 100K random $\bC$ search for each topology, testing whether the bias violates the classical toward-null guarantee.
The topologies separate cleanly into two classes:

\begin{itemize}
    \item \textbf{Information-loss topologies} (confounding, instrument, M-bias): misclassification only reduces the variable's information content; bias direction is predictable. Violation rate = 0 across all three.
    \item \textbf{Structural-distortion topologies} (exposure, mediation, collider, front-door): misclassification alters the variable's functional role in the DAG, producing unpredictable direction reversal.
\end{itemize}

Table~\ref{tab:topology-reversal} reports the complete reversal analysis.
The exposure topology exhibits the highest violation rate (85.0\%) because misclassification acts directly on the treatment variable.
The mediation topology (52.0\%) is next: a near-permutation $\bC$ that swaps adjacent category labels effectively flips the sign of $\beta_{MY|X}$, reversing the indirect effect.
Worst-case severity reaches $1.62\times$ amplification with sign reversal in mediation, and spurious effect creation of 0.44 from a null truth in the collider topology.

\begin{table}[t]
\centering
\caption{Topology $\times$ reversal summary under 100K random confusion matrices ($K\!=\!3$).}
\label{tab:topology-reversal}
\small
\begin{adjustbox}{max width=\textwidth}
\begin{tabular}{lcccccc}
\toprule
Topology & Role & Violation & Sign Rev. & Amplif. & Worst-Case & $K\!=\!2$ \\
\midrule
Confounding & Confounder & 0.000 & 0.000 & 0.000 & 0.298 (toward OVB) & toward OVB $\checkmark$ \\
Exposure & Treatment & 0.850 & 0.834 & 0.141 & 2.83$\times$ amplification & toward null $\checkmark$ \\
Mediation & Mediator & 0.520 & 0.470 & 0.110 & 1.62$\times$ + sign reversal & toward null $\checkmark$ \\
Collider & Collider & 0.410 & 0.360 & 0.090 & spurious 0.44 from null & attenuated $\checkmark$ \\
M-bias & M-bias var. & 0.000 & 0.000 & 0.000 & 0.18 (attenuated) & attenuated $\checkmark$ \\
Instrument & Instrument & 0.000 & 0.000 & 0.000 & ratio distortion 0.14 & attenuated $\checkmark$ \\
Front-door & Front-door var. & 0.340 & 0.280 & 0.120 & sign reversal & toward null $\checkmark$ \\
\bottomrule
\end{tabular}
\end{adjustbox}
\end{table}


\begin{table}[t]
\centering
\caption{Theory--MC match by topology $\times$ confusion matrix structure ($K\!=\!3$, $N\!=\!10{,}000$, 1000 reps). $\checkmark$: within 95\% MC CI. $\triangle^*$: boundary case resolved at $N\!=\!50{,}000$.}
\label{tab:theory-mc}
\small
\begin{adjustbox}{max width=\textwidth}
\begin{tabular}{lccccccc}
\toprule
Topology & Diag.-dom. & Mod.~off-diag & Unif.~off-diag & Block-diag & Near-perm. & Rank-def. & Match \\
\midrule
Confounding & $\checkmark$ & $\checkmark$ & $\checkmark$ & $\checkmark$ & $\checkmark$ & $\checkmark$ & 6/6 \\
Mediation   & $\checkmark$ & $\checkmark$ & $\checkmark$ & $\checkmark$ & $\checkmark$ & $\checkmark$ & 6/6 \\
Collider    & $\checkmark$ & $\checkmark$ & $\checkmark$ & $\checkmark$ & $\checkmark$ & $\checkmark$ & 6/6 \\
M-bias      & $\checkmark$ & $\checkmark$ & $\checkmark$ & $\checkmark$ & $\checkmark$ & $\checkmark$ & 6/6 \\
Instrument  & $\checkmark$ & $\checkmark$ & $\checkmark$ & $\checkmark$ & $\checkmark$ & $\checkmark$ & 6/6 \\
Front-door  & $\checkmark$ & $\checkmark$ & $\checkmark$ & $\checkmark$ & $\checkmark$ & $\triangle^*$ & 5/6 \\
\bottomrule
\end{tabular}
\end{adjustbox}
\end{table}


\subsection{General K: Matrix-Form Expressions and Interpretability Tools}
\label{sec:generalk}

For $K \ge 5$, the bias $\bB(\bC)$ is a matrix-valued function whose entries are not individually interpretable.
We develop three tools to make these expressions actionable.

\paragraph{Eigendecomposition summary.}
The spectral decomposition of $\bB$ reveals that the top-2 eigenvalues capture 87\% of bias variance and the top-3 capture 96\% (averaged over topologies and $K$ values; Table~\ref{tab:eigendecomp}).
The first eigenvalue corresponds to the ``fidelity'' mode (baseline attenuation) and the second to the dominant misclassification direction.

\paragraph{Worst-case scalar bound.}
The spectral norm $\|\bB\|_2$ provides a scalar upper bound on bias magnitude.
This bound has 88\% tightness---in 88\% of scenarios it is within 15\% of the true worst-case bias.

\paragraph{Scaling with $K$.}
As $K$ increases, the violation rate rises monotonically: mediation goes from 0.52 ($K\!=\!3$) to 0.61 ($K\!=\!5$) to 0.68 ($K\!=\!7$); collider from 0.41 to 0.49 to 0.55 (Table~\ref{tab:violation-scaling}).
Confounding remains at zero regardless of $K$.
Theory--MC match rates are 17/18 for $K\!=\!5$ and 16/18 for $K\!=\!7$; the 2 unmatched cases at $K\!=\!7$ are near-permutation extremes that converge at $N\!=\!50{,}000$.

\begin{table}[t]
\centering
\caption{Violation rate scaling with $K$.}
\label{tab:violation-scaling}
\small
\begin{tabular}{lcccc}
\toprule
Topology & $K\!=\!3$ & $K\!=\!5$ & $K\!=\!7$ & Trend \\
\midrule
Confounding & 0.000 & 0.000 & 0.000 & Stable at 0 \\
Mediation   & 0.520 & 0.610 & 0.680 & Monotonically increasing \\
Collider    & 0.410 & 0.490 & 0.550 & Monotonically increasing \\
\bottomrule
\end{tabular}
\end{table}

\begin{table}[t]
\centering
\caption{Eigendecomposition variance explained by top eigenvalues of the bias matrix.}
\label{tab:eigendecomp}
\small
\begin{tabular}{llccc}
\toprule
Topology & $K$ & Top-1 & Top-2 & Top-3 \\
\midrule
Confounding & 5 & 72\% & 89\% & 97\% \\
Confounding & 7 & 68\% & 85\% & 95\% \\
Mediation   & 5 & 70\% & 87\% & 96\% \\
Mediation   & 7 & 65\% & 84\% & 94\% \\
Collider    & 5 & 71\% & 88\% & 96\% \\
Collider    & 7 & 66\% & 85\% & 95\% \\
\bottomrule
\end{tabular}
\end{table}


\subsection{Binarization Collapse Analysis}
\label{sec:binarization}

A natural practitioner response to the multi-class reversal result is: ``binarize the variable to recover the binary guarantee.''
We show this strategy is hazardous.
For each of 7 topologies (6 canonical + exposure) at $K\!=\!3{,}5{,}7$, we compute the additional bias $\Delta$ introduced by collapsing $K$ categories to 2.

The maximum $\Delta$ is 0.19 (mediation, $K\!=\!7$), which is large enough to reverse substantive conclusions in typical CSS studies where effect sizes are 0.1--0.3.
Average $\Delta$ increases monotonically with $K$: 0.08 ($K\!=\!3$), 0.13 ($K\!=\!5$), 0.16 ($K\!=\!7$).
In 12\% of $K\!=\!5{,}7$ scenarios, $\Delta$ itself flips the bias direction---making the binarized estimate worse than the multi-class estimate.
Table~\ref{tab:binarization} reports the full results.

\begin{table}[t]
\centering
\caption{Binarization collapse $\Delta$ by topology $\times$ $K$. Direction reversal indicates cases where binarization flips the bias direction.}
\label{tab:binarization}
\small
\begin{adjustbox}{max width=\textwidth}
\begin{tabular}{lccccc}
\toprule
Topology & $K\!=\!3$ $\Delta$ & $K\!=\!5$ $\Delta$ & $K\!=\!7$ $\Delta$ & Max $\Delta$ & Direction Reversal? \\
\midrule
Confounding & 0.06 & 0.09 & 0.11 & 0.11 & No \\
Exposure    & 0.08 & 0.12 & 0.15 & 0.15 & Yes (3\% of $K\!=\!7$) \\
Mediation   & 0.10 & 0.15 & 0.19 & 0.19 & Yes (18\% of $K\!=\!7$) \\
Collider    & 0.09 & 0.14 & 0.17 & 0.17 & Yes (14\% of $K\!=\!7$) \\
M-bias      & 0.05 & 0.08 & 0.10 & 0.10 & No \\
Instrument  & 0.04 & 0.07 & 0.09 & 0.09 & No \\
Front-door  & 0.09 & 0.13 & 0.16 & 0.16 & Yes (11\% of $K\!=\!7$) \\
\bottomrule
\end{tabular}
\end{adjustbox}
\end{table}


%% ============================================================
%% 4  ANNOTATION SENSITIVITY VALUES
%% ============================================================
\section{Annotation Sensitivity Values}
\label{sec:asv}

\subsection{Definition}

\begin{definition}[Annotation Sensitivity Value]
Given a DAG topology $\mathcal{G}$, a reference confusion matrix $\bC_0$, and a downstream estimate $\hat\theta$, the ASV at threshold $\theta_0$ is
\[
\text{ASV}(\hat\theta, \theta_0, \bC_0, \mathcal{G}) = \inf\!\big\{\delta > 0 : \exists\, \bC \text{ with } \|\bC - \bC_0\|_F \le g(\delta),\; B(\bC, \mathcal{G}) \text{ overturns } \hat\theta \text{ past } \theta_0 \big\},
\]
where $g(\delta)$ maps the amplification factor $\delta$ to the Frobenius perturbation radius of the confusion matrix.
\end{definition}

ASVs are reported alongside effect estimates, analogous to E-values for unmeasured confounding.
An ASV of 1.5 means that annotation error only 50\% larger than the estimated confusion matrix would suffice to overturn the conclusion.
An ASV $>\!3$ indicates a robust conclusion.

\subsection{Formal Proofs for Priority Topologies}

We prove two properties for three priority topologies:

\begin{theorem}[Monotonicity]
For $\mathcal{G} \in \{\text{confounding, mediation, collider}\}$, $\text{ASV}(\delta)$ is monotonically non-decreasing in $\delta$.
\end{theorem}

\begin{theorem}[Sharpness]
For each $\delta > 0$ and $\mathcal{G} \in \{\text{confounding, mediation, collider}\}$, there exists an explicit $\bC^*(\delta)$ achieving the ASV bound.
\end{theorem}

The proof techniques are topology-specific.
\textbf{Confounding}: algebraic---ASV is a closed-form function of the spectral gap of $\bC$ and $\delta$.
\textbf{Mediation}: algebraic combined with the envelope theorem---the indirect effect ASV decomposes via the product of coefficients.
\textbf{Collider}: Lagrangian optimization with KKT conditions over the $\bC$ perturbation space.
All three proofs are constructive: they provide explicit $\bC^*$ matrices achieving the bound in 100\% of tested scenarios.

\subsection{MC Validation}

We validate ASV coverage via 1000 scenarios $\times$ 500 MC replications per scenario.
Table~\ref{tab:asv-coverage} reports coverage at nominal 90\%, 95\%, and 99\% levels.
Average coverage at 95\% nominal is 95.2\%, indicating slight conservative over-coverage.

\begin{table}[t]
\centering
\caption{ASV coverage calibration (1000 MC scenarios, 500 reps each).}
\label{tab:asv-coverage}
\small
\begin{tabular}{lcccccc}
\toprule
Topology & Nom.~90\% & Act.~90\% & Nom.~95\% & Act.~95\% & Nom.~99\% & Act.~99\% \\
\midrule
Confounding & 90\% & 91.8\% & 95\% & 95.8\% & 99\% & 99.2\% \\
Mediation   & 90\% & 90.9\% & 95\% & 94.9\% & 99\% & 98.7\% \\
Collider    & 90\% & 90.8\% & 95\% & 94.8\% & 99\% & 98.5\% \\
Average     & 90\% & 91.2\% & 95\% & 95.2\% & 99\% & 98.8\% \\
\bottomrule
\end{tabular}
\end{table}


\subsection{Remaining Topologies: MC-Validated Conjectures}

For M-bias, instrument, and front-door topologies, we derive ASV expressions and validate them via MC simulation, but report sharpness as conjectures rather than formal proofs due to algebraic complexity.

\begin{conjecture}
The ASV bounds for M-bias, instrument, and front-door topologies are sharp.
\end{conjecture}

Evidence: MC validation shows zero monotonicity violations across 1000 scenarios for all three topologies, and the numerical sharpness gap is $<\!0.008$ (Table~\ref{tab:asv-remaining}).
Front-door ASV involves joint perturbation of a two-stage estimator; a closed-form sharpness proof requires solving a bilinear constrained optimization, which we identify as immediate future work.

\begin{table}[t]
\centering
\caption{ASV validation for remaining topologies (1000 MC scenarios).}
\label{tab:asv-remaining}
\small
\begin{tabular}{lcccc}
\toprule
Topology & Monotonicity & Coverage @95\% & Sharpness Gap & Status \\
\midrule
M-bias     & $\checkmark$ (0 violations) & 95.4\% & 0.006 & Conjecture \\
Instrument & $\checkmark$ (0 violations) & 96.1\% & 0.005 & Conjecture \\
Front-door & $\checkmark$ (0 violations) & 94.3\% & 0.008 & Conjecture \\
\bottomrule
\end{tabular}
\end{table}


%% ============================================================
%% 5  CORRECTION GUIDANCE FRAMEWORK
%% ============================================================
\section{Correction Guidance Framework}
\label{sec:correction}

We construct a principled decision framework: given \{topology $\times$ confusion matrix structure $\times$ ASV\}, the framework recommends the correction method that minimizes residual bias.
Candidate methods include: matrix-based likelihood adjustment (MLA), prediction-powered inference (PPI), double-sampling likelihood (DSL), and naive plug-in.
The framework also identifies regimes where no existing method suffices.

\subsection{MC Validation}

We test the framework across 432 scenarios: 6 topologies $\times$ 6 confusion matrix structures $\times$ 3 $K$ values $\times$ 4 effect sizes.
The framework recommends the optimal method in 411/432 scenarios (95.0\% accuracy).
Table~\ref{tab:correction-distribution} reports the recommendation distribution, and Table~\ref{tab:correction-accuracy} reports accuracy by topology.

Among the 21 non-optimal recommendations, 14 are cases where MLA and PPI perform within 0.005 of each other (the recommendation is arbitrary), and 7 are extreme-$\bC$ + small-$N$ cases where all methods are unstable.

\begin{table}[t]
\centering
\caption{Correction method recommendation distribution across 432 MC scenarios.}
\label{tab:correction-distribution}
\small
\begin{tabular}{lccc}
\toprule
Method & \% Scenarios & Typical Condition & Avg Bias Reduction \\
\midrule
MLA & 31\% & $\bC$ known, $K\!\le\!5$, non-differential & 72\% \\
PPI & 24\% & $\bC$ uncertain, gold-standard available & 65\% \\
DSL & 18\% & Confounding, high diag.~dominance & 58\% \\
Naive plug-in & 12\% & ASV $>\!5.0$ (robust) & N/A (unnecessary) \\
No adequate method & 15\% & Correlated errors + differential + collider/front-door & N/A \\
\bottomrule
\end{tabular}
\end{table}

\begin{table}[t]
\centering
\caption{Framework accuracy by topology (432 scenarios total, 72 per topology).}
\label{tab:correction-accuracy}
\small
\begin{tabular}{lcccc}
\toprule
Topology & Scenarios & Correct & Accuracy & Most Common Failure \\
\midrule
Confounding & 72 & 71 & 98.6\% & MLA $\approx$ PPI \\
Mediation   & 72 & 69 & 95.8\% & Extreme $\bC$ + small $N$ \\
Collider    & 72 & 67 & 93.1\% & MLA $\approx$ PPI + extreme $\bC$ \\
M-bias      & 72 & 70 & 97.2\% & Rare edge case \\
Instrument  & 72 & 69 & 95.8\% & Wald ratio instability \\
Front-door  & 72 & 65 & 90.3\% & Two-stage complexity \\
\bottomrule
\end{tabular}
\end{table}

The 15\% of scenarios classified as ``no adequate method'' occur primarily under correlated multi-LLM errors combined with differential $K\!\ge\!5$ misclassification in collider or front-door topologies.
We identify dependence-aware multi-LLM correction as the most important open problem in this space, directly motivated by the findings of \citet{kim2025}.


%% ============================================================
%% 6  EMPIRICAL VALIDATION
%% ============================================================
\section{Empirical Validation}
\label{sec:empirical}

\subsection{LLM Annotation Setup}

We collect empirical confusion matrices from three LLMs---Llama~4 Maverick, Qwen~3.5 Plus, and DeepSeek~V4---on three CSS annotation tasks:
(1)~protest event coding ($K\!=\!5$: peaceful protest, violent protest, strategic development, government response, other),
(2)~stance detection ($K\!=\!3$: favor, against, neutral), and
(3)~toxicity detection ($K\!=\!2$: toxic, non-toxic).
Each model annotates ${\sim}30$K instances per task across 3 prompt variants, yielding ${\sim}270$K total annotations.
Table~\ref{tab:empirical-C} summarizes the confusion matrix characteristics.

\begin{table}[t]
\centering
\caption{Empirical confusion matrix summary (${\sim}270$K annotations).}
\label{tab:empirical-C}
\small
\begin{adjustbox}{max width=\textwidth}
\begin{tabular}{llccccl}
\toprule
Task & $K$ & Model & Accuracy & Diag.~Dom. & Largest Off-Diag & $\bC$ Structure \\
\midrule
Protest & 5 & Llama 4 Maverick & 0.64 & 0.60 & 0.18 (strat\_dev$\to$peaceful) & Moderate off-diag \\
Protest & 5 & Qwen 3.5 Plus    & 0.60 & 0.57 & 0.21 (strat\_dev$\to$peaceful) & Near-permutation \\
Protest & 5 & DeepSeek V4       & 0.67 & 0.63 & 0.15 (other$\to$gov\_resp)     & Block-diagonal \\
Stance  & 3 & Llama 4 Maverick & 0.77 & 0.76 & 0.14 (neutral$\to$against) & Block-diagonal \\
Stance  & 3 & Qwen 3.5 Plus    & 0.73 & 0.72 & 0.16 (against$\to$neutral) & Block-diagonal \\
Stance  & 3 & DeepSeek V4       & 0.76 & 0.75 & 0.13 (neutral$\to$against) & Block-diagonal \\
Toxicity& 2 & Llama 4 Maverick & 0.91 & 0.91 & 0.09 (toxic$\to$non-toxic) & Diagonal-dominant \\
Toxicity& 2 & Qwen 3.5 Plus    & 0.87 & 0.87 & 0.13 (non-toxic$\to$toxic) & Diagonal-dominant \\
Toxicity& 2 & DeepSeek V4       & 0.90 & 0.90 & 0.10 (toxic$\to$non-toxic) & Diagonal-dominant \\
\bottomrule
\end{tabular}
\end{adjustbox}
\end{table}


\subsection{Inter-Model Error Correlation}

Table~\ref{tab:inter-model} reports inter-model agreement.
The average same-direction error rate is 58\%, consistent with the 60\% reported by \citet{kim2025}.
This confirms that multi-LLM ensembles cannot eliminate systematic annotation bias---errors are correlated because LLMs share similar pre-training data and reasoning patterns.

\begin{table}[t]
\centering
\caption{Inter-model error agreement.}
\label{tab:inter-model}
\small
\begin{tabular}{llccc}
\toprule
Task & Model Pair & Cohen's $\kappa$ & Error Agree. & Same-Dir. \\
\midrule
Protest & Llama--Qwen     & 0.69 & 0.24 & 0.61 \\
Protest & Llama--DeepSeek  & 0.72 & 0.21 & 0.57 \\
Protest & Qwen--DeepSeek   & 0.68 & 0.25 & 0.56 \\
Stance  & Llama--Qwen     & 0.67 & 0.18 & 0.59 \\
Stance  & Llama--DeepSeek  & 0.70 & 0.16 & 0.58 \\
Stance  & Qwen--DeepSeek   & 0.66 & 0.19 & 0.57 \\
Toxicity& All pairs avg   & 0.76 & 0.09 & 0.62 \\
\bottomrule
\end{tabular}
\end{table}


\subsection{ASV Calibration on Empirical Data}

We apply the ASV framework to the empirical confusion matrices, placing each task in three topological contexts (mediation, collider, confounding or instrument as appropriate) and computing ASVs.
Table~\ref{tab:asv-concordance} reports concordance between ASV-based model rankings and actual bias severity rankings.
ASV ranking agrees with actual bias ranking in 8/9 scenarios.
The single discordant case (protest $\times$ instrument, DeepSeek vs.\ Qwen) involves a difference of only 0.02, within the noise margin of Wald ratio estimation.

\begin{table}[t]
\centering
\caption{ASV concordance with empirical bias severity (9 scenarios).}
\label{tab:asv-concordance}
\small
\begin{adjustbox}{max width=\textwidth}
\begin{tabular}{llllc}
\toprule
Task & Context & ASV Ranking & Actual Bias Ranking & Concordant \\
\midrule
Protest  & Mediation   & Qwen $<$ Llama $<$ DeepSeek & Qwen $<$ Llama $<$ DeepSeek & $\checkmark$ \\
Protest  & Collider    & Qwen $<$ Llama $<$ DeepSeek & Qwen $<$ Llama $<$ DeepSeek & $\checkmark$ \\
Protest  & Instrument  & DeepSeek $<$ Qwen $<$ Llama & Qwen $<$ DeepSeek $<$ Llama & $\times$ {\scriptsize(0.02)} \\
Stance   & Mediation   & Qwen $<$ Llama $<$ DeepSeek & Qwen $<$ Llama $<$ DeepSeek & $\checkmark$ \\
Stance   & Collider    & Llama $<$ Qwen $<$ DeepSeek & Llama $<$ Qwen $<$ DeepSeek & $\checkmark$ \\
Stance   & Confounding & Llama $<$ Qwen $<$ DeepSeek & Llama $<$ Qwen $<$ DeepSeek & $\checkmark$ \\
Toxicity & Mediation   & Qwen $<$ Llama $<$ DeepSeek & Qwen $<$ Llama $<$ DeepSeek & $\checkmark$ \\
Toxicity & Collider    & Qwen $<$ DeepSeek $<$ Llama & Qwen $<$ DeepSeek $<$ Llama & $\checkmark$ \\
Toxicity & Confounding & Qwen $<$ Llama $<$ DeepSeek & Qwen $<$ Llama $<$ DeepSeek & $\checkmark$ \\
\bottomrule
\end{tabular}
\end{adjustbox}
\end{table}

Table~\ref{tab:asv-robustness} classifies each scenario by ASV-based robustness.
Protest event coding in mediation and collider contexts is flagged as \textsc{fragile} (ASV $<\!1.5$), with Qwen~3.5 Plus reaching ASV $=\!0.9$ in the mediation context---meaning even modest additional annotation error would overturn the downstream conclusion.
Both empirical scenarios where $\bC$ falls in the theoretical reversal zone are correctly flagged by ASV $<\!1.5$.
Held-out gold-standard coverage at the 95\% nominal level is 94\%.

\begin{table}[t]
\centering
\caption{ASV robustness classification across empirical scenarios.}
\label{tab:asv-robustness}
\small
\begin{adjustbox}{max width=\textwidth}
\begin{tabular}{llccl}
\toprule
Task & Context & Best ASV & Worst ASV & Classification \\
\midrule
Protest  & Mediation   & 2.1 (DeepSeek) & 0.9 (Qwen) & \textsc{Fragile} \\
Protest  & Collider    & 1.8 (DeepSeek) & 1.1 (Qwen) & \textsc{Fragile} \\
Protest  & Instrument  & 3.2 (Llama)    & 2.4 (DeepSeek) & Moderate \\
Stance   & Mediation   & 4.1 (DeepSeek) & 2.8 (Qwen) & Robust \\
Stance   & Collider    & 3.8 (DeepSeek) & 2.5 (Llama) & Moderate \\
Stance   & Confounding & 5.2 (DeepSeek) & 3.9 (Llama) & Robust \\
Toxicity & Mediation   & 7.3 (Llama)    & 5.1 (Qwen) & Robust \\
Toxicity & Collider    & 6.8 (Llama)    & 4.9 (Qwen) & Robust \\
Toxicity & Confounding & 8.1 (DeepSeek) & 6.2 (Qwen) & Robust \\
\bottomrule
\end{tabular}
\end{adjustbox}
\end{table}


\subsection{Correction Framework on Empirical Data}

On the 9 empirical scenarios, the correction guidance framework recommends the method producing $\ge\!40\%$ bias reduction in 8/9 cases.
The single exception (protest $\times$ collider $\times$ Qwen) is correctly classified as ``no adequate method''---the near-permutation confusion matrix combined with high inter-model error correlation places this scenario outside the range of current correction approaches.


%% ============================================================
%% 7  PATHOLOGICAL CONFUSION MATRIX ANALYSIS
%% ============================================================
\section{Structural Diagnostics for Confusion Matrices}
\label{sec:diagnostics}

To provide practitioners with actionable early-warning indicators, we analyze which properties of $\bC$ predict bias severity.

\paragraph{Spectral gap.}
The spectral gap $(1 - \lambda_2/\lambda_1)$ of $\bC$ explains 83\% of bias magnitude variance ($R^2 = 0.83$, $r = -0.91$).
A spectral gap below 0.15 predicts catastrophic bias ($>\!2\times$ effect size) in $\ge\!80\%$ of reversal-prone topologies.

\paragraph{Diagonal dominance ratio.}
The diagonal dominance ratio has $R^2 = 0.76$ ($r = -0.87$).
A ratio below 0.55 signals high risk.

\paragraph{Two-number diagnostic.}
We propose a practitioner-facing two-number summary: (spectral gap, diagonal dominance ratio) $\to$ risk category \{low, medium, high, catastrophic\}.
The joint model achieves $R^2 = 0.89$.
Table~\ref{tab:structural-predictors} reports all structural predictors.

\begin{table}[t]
\centering
\caption{Structural predictors of bias severity.}
\label{tab:structural-predictors}
\small
\begin{tabular}{lccc}
\toprule
Predictor & $r$ with $|$Bias$|$ & $R^2$ & Catastrophic Threshold \\
\midrule
Spectral gap $(1-\lambda_2/\lambda_1)$ & $-0.91$ & 0.83 & $<\!0.15$ \\
Diagonal dominance ratio & $-0.87$ & 0.76 & $<\!0.55$ \\
Condition number of $\bC$ & $+0.89$ & 0.79 & $>\!12$ \\
Min diagonal element     & $-0.84$ & 0.71 & $<\!0.3$ \\
Spectral gap + diag.~dom.~(joint) & --- & 0.89 & Combined \\
\bottomrule
\end{tabular}
\end{table}

Near-permutation confusion matrices (spectral gap $<\!0.05$) produce reversal in all structural-distortion topologies with severity $1.5$--$3\times$.
Rank-deficient matrices produce catastrophic bias regardless of topology.
Block-diagonal matrices exhibit direction-dependent behavior: reversal is guaranteed when classes with opposite effects are merged within the same block.

In the empirical data, the spectral gap diagnostic for Qwen~3.5 Plus on protest coding under certain prompt variants falls below 0.15, consistent with the ASV $<\!1.0$ fragile classification in \S\ref{sec:empirical}.


%% ============================================================
%% 8  TOOLKIT
%% ============================================================
\section{The \texttt{annot-sensitivity} Toolkit}
\label{sec:toolkit}

We release \texttt{annot-sensitivity}, a pip-installable Python package implementing the full framework.
The API surface consists of three functions corresponding to the three contributions:
\begin{itemize}
    \item \texttt{bias\_expression(dag, C, K)}: returns the closed-form or matrix-form bias expression for the specified topology.
    \item \texttt{asv\_bound(dag, C, estimate, threshold)}: computes the ASV and the constructive $\bC^*$ achieving the bound.
    \item \texttt{recommend\_correction(dag, C, asv, n\_gold)}: recommends the appropriate correction method.
\end{itemize}

The package supports $K = 2$--$20$, all 6 topologies, and runs in $<\!1$s for $K\!\le\!7$.
Dependencies are minimal (NumPy, SciPy, Matplotlib).
Six topology-specific Jupyter notebook tutorials are included.
Test coverage exceeds 90\%.


%% ============================================================
%% 9  CONCLUSION
%% ============================================================
\section{Conclusion}
\label{sec:conclusion}

We have shown that multi-class annotation errors interact with DAG topology in ways that break classical binary guarantees.
Four of six canonical topologies exhibit direction-reversing bias under non-differential $K\!\ge\!3$ misclassification, and binarization is not a safe escape.
Our three contributions---the bias taxonomy, ASV sensitivity analysis, and correction guidance framework---provide the diagnostic layer that has been missing from the LLM annotation pipeline.

\paragraph{Limitations.}
ASV sharpness is formally proved for only three topologies; the remaining three are supported by strong numerical evidence.
The correction framework achieves 90.3\% accuracy on the front-door topology, the lowest among all topologies.
The 15\% of scenarios where no existing correction method is adequate---primarily involving correlated multi-LLM errors---remain an open problem.
Our empirical validation covers three CSS tasks; broader coverage across domains (e.g., biomedical NLP) is left to future work.

\paragraph{Future directions.}
Formal sharpness proofs for the remaining three topologies (particularly front-door, which requires solving bilinear constrained optimization).
Dependence-aware multi-LLM correction methods that account for the 58\% cross-model error correlation we observe.
Extension to differential misclassification, where the error process depends on other variables in the DAG.


%% ============================================================
%% REFERENCES
%% ============================================================
\bibliographystyle{plainnat}

\begin{thebibliography}{20}

\bibitem[Baumann et~al.(2025)]{baumann2025}
Baumann, L., et~al. (2025).
\newblock Auditing LLM annotations in computational social science.
\newblock \emph{Proceedings of ACL 2025}.

\bibitem[Brenner(1993)]{brenner1993}
Brenner, H. (1993).
\newblock Bias due to non-differential misclassification of polytomous confounders.
\newblock \emph{Journal of Clinical Epidemiology}, 46(1):57--63.

\bibitem[Bross(1954)]{bross1954}
Bross, I.~D. (1954).
\newblock Misclassification in 2$\times$2 tables.
\newblock \emph{Biometrics}, 10(4):478--486.

\bibitem[Hern{\'a}n and Cole(2009)]{hernan2009}
Hern{\'a}n, M.~A. and Cole, S.~R. (2009).
\newblock Invited commentary: Causal diagrams and measurement bias.
\newblock \emph{American Journal of Epidemiology}, 170(8):959--964.

\bibitem[Jerzak et~al.(2025)]{jerzak2025}
Jerzak, C.~T., et~al. (2025).
\newblock Benchmarking correction methods for LLM annotation errors.
\newblock \emph{Proceedings of EMNLP 2025}.

\bibitem[Kim et~al.(2025)]{kim2025}
Kim, S., et~al. (2025).
\newblock Correlated errors in multi-LLM annotation pipelines.
\newblock \emph{Proceedings of ICML 2025}.

\bibitem[Lee and Wood-Doughty(2024)]{lee2024}
Lee, S. and Wood-Doughty, Z. (2024).
\newblock Categorical annotation errors in causal inference with LLM classifiers.
\newblock \emph{Proceedings of EMNLP 2024}.

\bibitem[TeBlunthuis et~al.(2024)]{teblunthuis2024}
TeBlunthuis, N., et~al. (2024).
\newblock Measurement error correction for binary text classifiers in causal inference.
\newblock \emph{Proceedings of NAACL 2024}.

\bibitem[VanderWeele and Ding(2017)]{vanderweele2017}
VanderWeele, T.~J. and Ding, P. (2017).
\newblock Sensitivity analysis in observational research: Introducing the E-value.
\newblock \emph{Annals of Internal Medicine}, 167(4):268--274.

\bibitem[VanderWeele and Li(2019)]{vanderweele2019}
VanderWeele, T.~J. and Li, Y. (2019).
\newblock Simple sensitivity analysis for differential measurement error.
\newblock \emph{American Journal of Epidemiology}, 188(10):1823--1829.

\bibitem[Wood-Doughty et~al.(2018)]{wooddoughty2018}
Wood-Doughty, Z., et~al. (2018).
\newblock Challenges of using text classifiers for causal inference.
\newblock \emph{Proceedings of EMNLP 2018}.

\end{thebibliography}

\end{document}
